% !TeX root = ../../../main.tex
\subsubsection{ITE Distribution for Different Medication}
\label{ch03:sec:04_results_ite_analysis}

We conducted a subgroup analysis to compare the distribution of Estimated Individual Treatment Effects (ITE) across three adherence groups: patients covered by injectable medications, non-injectable medications, and those not covered by any medication at the time snapshot \(\tau=3\) months.
Figure~\ref{ch03:fig:ite_injectable} presents the ITE distributions for each adherence group using various survival models and meta-learners.
The ITEs reflect the change in Restricted Mean Event Time (RMET) when comparing a non-adherence scenario to adherence within each category.

For the T-learner-based approaches using survival models, the ATEs for injectable medications appear slightly more negative compared to the non-injectable or not-covered groups.
This trend is observed in Figures~\ref{ch03:fig:ite_injectable_coxph}, \ref{ch03:fig:ite_injectable_deepsurv}, and \ref{ch03:fig:ite_injectable_random_survival_forest}, corresponding to CoxPH, DeepSurv, and Random Survival Forest, respectively.
However, the Causal Survival Forest results in Figure~\ref{ch03:fig:ite_injectable_causal_survival_forest} do not show a similar distinction, with ATEs for injectables being closer to the non-injectable and not-covered groups.

Overall, the ATEs across all survival models remain close within each formulation group, suggesting no meaningful differences in RMET across injectable, non-injectable, and not-covered groups in our dataset.
The ITE distributions also highlight some variability across models, with Random Survival Forest showing relatively narrower distributions compared to the other methods.
Further subgroup analysis is detailed in Appendix~\ref{ch03:apd:ite_injectable_all} for all models and meta-learners and follow the same pattern observed in Figure~\ref{ch03:fig:ite_injectable}.

These consistent negative ATE values align with the results reported in Section~\ref{ch03:sec:04_results_causal}, indicating shorter adverse event times associated with non-adherence across all adherence categories and medication types.
If nuanced differences existed in ITEs for injectable versus non-injectable formulations, they might have highlighted heterogeneity in treatment effects based on the delivery mechanism of medications.
However, the lack of such variation suggests that the adherence effect may be largely uniform across these groups in the context of our dataset.

In Section~\ref{ch03:apd:ite_medication}, we conducted a similar analysis for ITEs by medication type, observing close ATE values across the top four antipsychotic medications --risperidone, aripiprazole, olanzapine, and haloperidol-- with minimal differences in RMETs for non-adherence within each medication, showing no heterogeneity across medication type as well.

These findings indicate no significant heterogeneity in treatment effects across medication formulations or specific antipsychotic medications.
While slight variations in ATE rankings are observed across models, the magnitudes remain similar, suggesting that the adherence effect is largely stable regardless of medication type or delivery mechanism.

